\section{Threats to Validity}
\label{sec:threats}

\subsection{Internal Validity}

The main internal validity risk is optimistic bias from design or
hyperparameter choices. To reduce this risk, the WR-STV top-\(k\)
parameter and residual clip bound \(C\) are selected using a validation
split independent of the held-out test traces, as described in
Section~\ref{sec:hyperparameter-selection}. The primary injection
setting \((f=30, m=3)\) is reported together with injection
strength sensitivity over \(f \in \{5,10,20,30\}\), so the effect of injection
strength is visible rather than hidden in a single operating point.
All STV vocabularies, residual statistics, and reliability weights are
estimated from training traces only and then applied unchanged to
validation or test traces.

\subsection{Construct Validity}

The evaluation measures controlled synthetic latency anomalies rather
than naturally occurring production failures. Injected anomalies are
created by multiplying selected span durations, so the evaluation
primarily tests sensitivity to path-visible latency deviations. It does
not test non-latency anomalies such as incorrect responses, missing
calls, dropped spans, novel service paths, or structural changes in the
trace graph. In addition, both largest-visible and random-visible
injection strategies are constrained to STV-visible service paths, i.e.,
features already present in the training vocabulary. Therefore, the
results characterize WR-STV under controlled STV-visible latency
perturbations rather than all possible microservice anomaly types.

The reported F1 score is computed using an oracle threshold that labels
the top \(N_+\) scored traces as anomalous, where \(N_+\) is the number
of injected traces in the evaluation split. This metric should be
interpreted as a ranking diagnostic rather than as expected deployment
precision or recall under an unknown operational threshold.

\subsection{External Validity}

The evaluation is limited to three configuration variants of
\texttt{ts-order-service} within the TrainTicket benchmark. Although
these configurations provide useful version-induced variation, they
represent a single benchmark system and a single target service. The
results may not generalize to other services, other microservice
benchmarks, production workloads, or environments with substantially
different tracing quality, traffic patterns, or service-call structures.
Additional evaluation on other services, benchmark systems, and real
production incident data is required before making broader
generalization claims.

\subsection{Statistical Conclusion Validity}

Within-dataset statistical analyses are based on repeated experiments
over five random seeds for each of three datasets, yielding \(n=15\)
seed-dataset runs. These repeated runs improve robustness to random
train-test splitting, but they are not equivalent to fifteen independent
datasets. Cross-dataset transfer analysis is even more limited: it uses
only six directed train-test pairs derived from the same three datasets.
Consequently, cross-dataset statistical tests should be interpreted as
exploratory evidence supporting observed trends rather than definitive
statistical confirmation.